\documentclass[11pt]{article}
\newcommand{\documenttype}{Cover letter}
\usepackage[letterpaper,margin=1in]{geometry}
\usepackage[T1]{fontenc}
\usepackage[utf8]{inputenc}
\usepackage{newtxtext,newtxmath}
\usepackage{microtype}
\usepackage{xcolor}
\usepackage{enumitem}
\usepackage{tabularx}
\usepackage{titlesec}
\usepackage{fancyhdr}
\usepackage{hyperref}
\definecolor{accent}{HTML}{315B8A}
\definecolor{link}{HTML}{5B4F7A}
\hypersetup{colorlinks=true,urlcolor=link,linkcolor=link,citecolor=link,
  pdfauthor={Phurailatpam Hemantakumar}}
\setlength{\parindent}{0pt}
\setlength{\parskip}{6pt}
\setlength{\headheight}{14pt}
\setlength{\emergencystretch}{1.5em}
\titleformat{\section}{\large\bfseries\color{accent}}{}{0pt}{}
\titlespacing*{\section}{0pt}{11pt}{5pt}
\titleformat{\subsection}{\normalsize\bfseries}{}{0pt}{}
\titlespacing*{\subsection}{0pt}{8pt}{3pt}
\setlist[itemize]{leftmargin=15pt,itemsep=3pt,topsep=3pt,parsep=0pt}
\setlist[enumerate]{leftmargin=19pt,itemsep=7pt,topsep=4pt,parsep=0pt}
\pagestyle{fancy}
\fancyhf{}
\fancyhead[L]{\small \href{https://hemantaph.com}{Phurailatpam Hemantakumar}}
\fancyhead[R]{\small\documenttype}
\fancyfoot[R]{\small\thepage}
\renewcommand{\headrulewidth}{0.3pt}
\renewcommand{\footrulewidth}{0pt}
\newcommand{\firstheading}[2]{%
  \thispagestyle{plain}%
  {\LARGE\bfseries\color{accent} #1}\par
  \vspace{3pt}{#2}\par
  \vspace{4pt}{\color{accent}\hrule}\vspace{8pt}}
\newcommand{\dated}[2]{\par\noindent\begin{tabularx}{\linewidth}{@{}X r@{}}#1 & #2\end{tabularx}\par}

% Shared contact details for the NBI CV and cover letter.
\newcommand{\contact}{%
  \href{mailto:hemantaph@link.cuhk.edu.hk}{hemantaph@link.cuhk.edu.hk}\quad |\quad +852 9867 8178\\
  \href{https://hemantaph.com/}{hemantaph.com}\quad |\quad The Chinese University of Hong Kong}

\begin{document}
\firstheading{Phurailatpam Hemantakumar}{\contact}
13 November 2026

Postdoctoral Fellowship Search Committee\\
Niels Bohr International Academy, Niels Bohr Institute\\
University of Copenhagen

Dear Prof. Pessah and Members of the Search Committee,

I am applying for the Postdoctoral Fellowship in Theoretical and Gravitational Wave Astrophysics (ASTRO27PDR). I completed my PhD research in Physics at The Chinese University of Hong Kong under the supervision of Otto Akseli Hannuksela. The formal award of my degree is expected in December 2026. My research focuses on gravitational-wave (GW) lensing statistics, eccentric waveform development and implementation, and Bayesian data analysis. I also develop computational methods for population simulations and detection-probability calculations.

I lead the development of \href{https://ler.hemantaph.com/}{\texttt{ler}}, an open-source Monte Carlo framework for simulating lensed and unlensed compact-binary populations and predicting detectable event rates. I also developed \href{https://gwsnr.hemantaph.com/}{\texttt{gwsnr}} for efficient signal-to-noise ratio and detection-probability calculations. My work involves astrophysical population models, detector-selection effects, numerical methods, uncertainty propagation, and statistical analysis. I am currently working with Chunglee Kim and Jiyun Park on star-formation-informed binary-black-hole merger rates. I have also gained familiarity with binary population synthesis by reproducing results from \href{https://arxiv.org/abs/2405.01623}{Boesky et al. (2024)} using \href{https://compas.science/}{COMPAS} for a \href{https://docs.google.com/presentation/d/1WBn-nu949xdn9oGZe1GgIlbHsQ_ZBelahJA6aSFyVq0/edit}{journal club presentation}.

With collaborators, I co-developed \href{https://arxiv.org/abs/2508.12697}{\texttt{TaylorF2Ecck}}, an analytic, frequency-domain post-Newtonian eccentric inspiral waveform, implemented it in a \texttt{LALSuite} fork, and performed Bayesian analyses of GW170817 and GW190425. In a related project with Lalit Pathak and Achamveedu Gopakumar, I investigated the degeneracy between orbital eccentricity and line-of-sight acceleration induced by a possible tertiary companion in GW190814.

The programme aims to bring together astrophysical fluid dynamics and GW astrophysics to understand the environments leading to compact-binary mergers. My expertise in GW population simulations, statistical analysis, Bayesian data analysis, and analytical waveform modelling could complement this effort. I would welcome discussions with the group about contributions that build on these methods and my existing research.

Alongside contributing to the group's research programme, I intend to pursue independent research in GW lensing statistics, eccentric waveform modelling and data analysis, and multimessenger lensing, building on my existing projects and collaborations. My research statement describes these plans, including the continuation of our line-of-sight acceleration study. I will also continue my existing LVK/GRASP responsibilities, including leadership and development of \texttt{ler} towards GWTC-6 and GWTC-7 and continued development and support of \texttt{gwsnr}. I have mentored undergraduate and graduate researchers, taught scientific Python, and conducted GW data-analysis workshops, and would welcome opportunities to contribute to mentoring at NBIA.

Thank you for considering my application.

Sincerely,\\
Phurailatpam Hemantakumar
\end{document}
